% GENERATED FILE. Edit canonical lessons, then run npm run generate:books.
% canonical-source-hash: fnv1a64:9d60158b73524b60
% canonical-lessons: MR-C256-usne-ghene, MR-C256-japne, MR-C256-vyakta-karne, MR-C256-barik-hone, MR-C256-khau-ghalne

\chapter{To borrow, To look after, To express, To get thin, To feed}
\label{ch:mr-to-borrow-to-look-after-to-express-to-get-thin-to-feed}
\hlchaptermodality{256}

\begin{chapteropening}
\textbf{By the end of this chapter:} \emph{I can say to borrow, to look after, to express, to get thin and to feed.}

Complete the last lesson of chapter 256: I can say to borrow, to look after, to express, to get thin and to feed.
\end{chapteropening}

\section[\texorpdfstring{\mr{उसने} \mr{घेणे} (usne gheṇe)}{उसने घेणे (usne gheṇe)}]{\mr{उसने} \mr{घेणे} (usne gheṇe) — to borrow}

\label{lesson:MR-C256-usne-ghene}



\begin{quote}
Before the new one: say the Marathi for to pray, then the Marathi for to recommend.
\end{quote}

\subsection*{You'll want to know: \mr{उसने} \mr{घेणे}}
\textbf{\mr{उसने} \mr{घेणे}} — \emph{usne gheṇe} — \textquotedblleft{}to borrow\textquotedblright{}.

\begin{grammarlens}[title={What you've built}]
One more word for this chapter.
\end{grammarlens}

\subsection*{Your turn}
\begin{itemize}
\raggedright
  \item \emph{Say it:} \emph{usne gheṇe}
  \item \emph{Say it:} \emph{usne gheṇe}, once more
  \item \emph{Say it:} say \emph{usne gheṇe}
  \item \emph{Recall:} say the Marathi for to pray, then the Marathi for to recommend, then say \emph{usne gheṇe} again
\end{itemize}

\begin{tcolorbox}[breakable,colback=teal!4,colframe=teal!35!black,title={Before you move on}]
What does \mr{उसने} \mr{घेणे} mean? (\textquotedblleft{}To borrow\textquotedblright{}.) Say it once more.
\end{tcolorbox}
\section[\texorpdfstring{\mr{जपणे} (japṇe)}{जपणे (japṇe)}]{\mr{जपणे} (japṇe) — to look after, to cherish}

\label{lesson:MR-C256-japne}



\begin{quote}
Before the new one: say the Marathi for to recommend, then the Marathi for to borrow.
\end{quote}

\subsection*{You'll want to know: \mr{जपणे}}
\textbf{\mr{जपणे}} — \emph{japṇe} — \textquotedblleft{}to look after, to cherish\textquotedblright{}.

\begin{grammarlens}[title={What you've built}]
One more word for this chapter.
\end{grammarlens}

\subsection*{Your turn}
\begin{itemize}
\raggedright
  \item \emph{Say it:} \emph{japṇe}
  \item \emph{Say it:} \emph{japṇe}, once more
  \item \emph{Say it:} say \emph{japṇe}
  \item \emph{Recall:} say the Marathi for to recommend, then the Marathi for to borrow, then say \emph{japṇe} again
\end{itemize}

\begin{tcolorbox}[breakable,colback=teal!4,colframe=teal!35!black,title={Before you move on}]
What does \mr{जपणे} mean? (\textquotedblleft{}To look after, to cherish\textquotedblright{}.) Say it once more.
\end{tcolorbox}
\section[\texorpdfstring{\mr{व्यक्त} \mr{करणे} (vyakta karṇe)}{व्यक्त करणे (vyakta karṇe)}]{\mr{व्यक्त} \mr{करणे} (vyakta karṇe) — to express}

\label{lesson:MR-C256-vyakta-karne}



\begin{quote}
Before the new one: say the Marathi for to borrow, then the Marathi for to look after.
\end{quote}

\subsection*{You'll want to know: \mr{व्यक्त} \mr{करणे}}
\textbf{\mr{व्यक्त} \mr{करणे}} — \emph{vyakta karṇe} — \textquotedblleft{}to express\textquotedblright{}.

\begin{grammarlens}[title={What you've built}]
One more word for this chapter.
\end{grammarlens}

\subsection*{Your turn}
\begin{itemize}
\raggedright
  \item \emph{Say it:} \emph{vyakta karṇe}
  \item \emph{Say it:} \emph{vyakta karṇe}, once more
  \item \emph{Say it:} say \emph{vyakta karṇe}
  \item \emph{Recall:} say the Marathi for to borrow, then the Marathi for to look after, then say \emph{vyakta karṇe} again
\end{itemize}

\begin{tcolorbox}[breakable,colback=teal!4,colframe=teal!35!black,title={Before you move on}]
What does \mr{व्यक्त} \mr{करणे} mean? (\textquotedblleft{}To express\textquotedblright{}.) Say it once more.
\end{tcolorbox}
\section[\texorpdfstring{\mr{बारीक} \mr{होणे} (bārīk hoṇe)}{बारीक होणे (bārīk hoṇe)}]{\mr{बारीक} \mr{होणे} (bārīk hoṇe) — to get thin}

\label{lesson:MR-C256-barik-hone}



\begin{quote}
Before the new one: say the Marathi for to look after, then the Marathi for to express.
\end{quote}

\subsection*{You'll want to know: \mr{बारीक} \mr{होणे}}
\textbf{\mr{बारीक} \mr{होणे}} — \emph{bārīk hoṇe} — \textquotedblleft{}to get thin\textquotedblright{}.

\begin{grammarlens}[title={What you've built}]
One more word for this chapter.
\end{grammarlens}

\subsection*{Your turn}
\begin{itemize}
\raggedright
  \item \emph{Say it:} \emph{bārīk hoṇe}
  \item \emph{Say it:} \emph{bārīk hoṇe}, once more
  \item \emph{Say it:} say \emph{bārīk hoṇe}
  \item \emph{Recall:} say the Marathi for to look after, then the Marathi for to express, then say \emph{bārīk hoṇe} again
\end{itemize}

\begin{tcolorbox}[breakable,colback=teal!4,colframe=teal!35!black,title={Before you move on}]
What does \mr{बारीक} \mr{होणे} mean? (\textquotedblleft{}To get thin\textquotedblright{}.) Say it once more.
\end{tcolorbox}
\section[\texorpdfstring{\mr{खाऊ} \mr{घालणे} (khāū ghālṇe)}{खाऊ घालणे (khāū ghālṇe)}]{\mr{खाऊ} \mr{घालणे} (khāū ghālṇe) — to feed}

\label{lesson:MR-C256-khau-ghalne}



\begin{quote}
Before the new one: say the Marathi for to express, then the Marathi for to get thin.
\end{quote}

\subsection*{You'll want to know: \mr{खाऊ} \mr{घालणे}}
\textbf{\mr{खाऊ} \mr{घालणे}} — \emph{khāū ghālṇe} — \textquotedblleft{}to feed\textquotedblright{}.

\begin{grammarlens}[title={What you've built}]
One more word for this chapter.
\end{grammarlens}

\subsection*{Your turn}
\begin{itemize}
\raggedright
  \item \emph{Say it:} \emph{khāū ghālṇe}
  \item \emph{Say it:} \emph{khāū ghālṇe}, once more
  \item \emph{Say it:} say \emph{khāū ghālṇe}
  \item \emph{Recall:} say the Marathi for to express, then the Marathi for to get thin, then say \emph{khāū ghālṇe} again
\end{itemize}

\begin{tcolorbox}[breakable,colback=teal!4,colframe=teal!35!black,title={Before you move on}]
What does \mr{खाऊ} \mr{घालणे} mean? (\textquotedblleft{}To feed\textquotedblright{}.) Say it once more.
\end{tcolorbox}
